% Professional resume — pdfLaTeX, no custom template required
% Compile: make professional

\documentclass[11pt,letterpaper]{article}
\usepackage[margin=0.65in]{geometry}
\usepackage{enumitem}
\usepackage{titlesec}
\usepackage{hyperref}
\usepackage{tabularx}

\hypersetup{colorlinks=true, urlcolor=blue, linkcolor=black}
\pagestyle{empty}
\setlength{\parindent}{0pt}
\setlength{\parskip}{0pt}

\titleformat{\section}{\large\bfseries\uppercase}{}{0pt}{}[\titlerule]
\titlespacing*{\section}{0pt}{10pt}{4pt}

\newcommand{\role}[4]{%
  \textbf{#1} \hfill #2\\
  \textit{#3} \hfill \textit{#4}\\[-2pt]
}

\begin{document}

\begin{center}
  {\LARGE\bfseries Dr. Julia Hasty}\\[6pt]
  \href{mailto:jay@resilientmarkets.com}{jay@resilientmarkets.com} \textbullet{}
  \href{https://www.linkedin.com/in/dr-jay/}{linkedin.com/in/dr-jay} \textbullet{}
  540.798.2312 \textbullet{} Albany, CA
\end{center}

\vspace{4pt}
\begin{center}
\small
PhD materials scientist \textbullet{} XPrize winner \textbullet{}
900+ clean energy deployments in 35+ countries \textbullet{}
microgrids, water systems, energy policy
\end{center}

\section{Experience}

\role{Resilient Markets}{2020 -- Present}{Managing Partner}{Albany, CA}
\begin{itemize}[nosep, leftmargin=*]
  \item Technical assessments of biomass, solar, and battery microgrids for critical infrastructure, agriculture, and water
  \item Policy and regulatory research with MaxETA for the Puerto Rico Energy Bureau
\end{itemize}

\role{10Power}{2020 -- Present}{Technical Lead}{San Francisco, CA}
\begin{itemize}[nosep, leftmargin=*]
  \item Published the 2024 Energy Plan for the Alaska Native Village of Eklutna
  \item Grant writing and integrated microgrid development for Native American tribal communities
\end{itemize}

\role{Corvidae, LLC}{2017 -- 2020}{Founder}{Oakland, CA}
\begin{itemize}[nosep, leftmargin=*]
  \item Advised 15+ early-stage clean-tech companies; led Puerto Rico microgrid portfolio and co-founded ACEPR cooperatives accelerator
\end{itemize}

\role{SkySource}{2018}{Systems Engineer}{Berkeley, CA}
\begin{itemize}[nosep, leftmargin=*]
  \item \textbf{Winner, 2018 Water Abundance XPrize}: biomass-powered atmospheric water (2,000 L/day at $\sim$2\textcent/L)
\end{itemize}

\role{ALL Power Labs}{2009 -- 2017}{Engineering Director}{Berkeley, CA}
\begin{itemize}[nosep, leftmargin=*]
  \item Directed 15 engineers; 900+ deployments of 150 kW gasification systems in 35+ countries
  \item NSF Principal Investigator; USAID project lead (250 homes, Liberia)
\end{itemize}

\role{Earlier roles}{2007 -- 2021}{}{}
CPWater Tech, Box Power, NanoSulf, Portland Biodiesel; American Red Cross NAT Lab (FDA-regulated testing); Stony Brook PhD research and adjunct teaching.

\section{Education}

\textbf{Ph.D., Materials Science \& Engineering} (Chemical Engineering focus) --- Stony Brook University\\
\textbf{B.S., Chemistry} --- Radford University

\section{Selected publications \& patents}

Patents US 8,829,695 \& US 9,476,352 (compact gasifier-genset architecture, 2017).\\
\textit{Fuel} (2016); \textit{J. Renewable Sustainable Energy} (2015). Full list: \texttt{publications-supplement.pdf}.

\section{Skills}

\textbf{Laboratory:} NMR, FTIR, UV-VIS, HPLC, PCR, GC-MS, SEM, TEM\\
\textbf{Engineering:} Gasification, microgrids, water systems, automation\\
\textbf{Tools:} Python, \LaTeX, Solidworks, ChemCAD, Fusion360

\end{document}
